% The Slide Definitions
%document
\documentclass[10pt]{beamer}
%theme
\usetheme{metropolis}
% packages
\usepackage{color}
\usepackage{listings}
\usepackage[ngerman]{babel}
\usepackage[utf8]{inputenc}
\usepackage{multicol}


% color definitions
\definecolor{mygreen}{rgb}{0,0.6,0}
\definecolor{mygray}{rgb}{0.5,0.5,0.5}
\definecolor{mymauve}{rgb}{0.58,0,0.82}

\lstset{
    backgroundcolor=\color{white},
    % choose the background color;
    % you must add \usepackage{color} or \usepackage{xcolor}
    basicstyle=\footnotesize\ttfamily,
    % the size of the fonts that are used for the code
    breakatwhitespace=false,
    % sets if automatic breaks should only happen at whitespace
    breaklines=true,                 % sets automatic line breaking
    captionpos=b,                    % sets the caption-position to bottom
    commentstyle=\color{mygreen},    % comment style
    % deletekeywords={...},
    % if you want to delete keywords from the given language
    extendedchars=true,
    % lets you use non-ASCII characters;
    % for 8-bits encodings only, does not work with UTF-8
    frame=single,                    % adds a frame around the code
    keepspaces=true,
    % keeps spaces in text,
    % useful for keeping indentation of code
    % (possibly needs columns=flexible)
    keywordstyle=\color{blue},       % keyword style
    % morekeywords={*,...},
    % if you want to add more keywords to the set
    numbers=left,
    % where to put the line-numbers; possible values are (none, left, right)
    numbersep=5pt,
    % how far the line-numbers are from the code
    numberstyle=\tiny\color{mygray},
    % the style that is used for the line-numbers
    rulecolor=\color{black},
    % if not set, the frame-color may be changed on line-breaks
    % within not-black text (e.g. comments (green here))
    stepnumber=1,
    % the step between two line-numbers.
    % If it's 1, each line will be numbered
    stringstyle=\color{mymauve},     % string literal style
    tabsize=4,                       % sets default tabsize to 4 spaces
    % show the filename of files included with \lstinputlisting;
    % also try caption instead of title
    language = [Sharp]C,
	showspaces = false,
	showtabs = false,
	showstringspaces = false,
	escapechar = ,
}

\def\ContinueLineNumber{\lstset{firstnumber=last}}
\def\StartLineAt#1{\lstset{firstnumber=#1}}
\let\numberLineAt\StartLineAt



\newcommand{\codeline}[1]{
	\alert{\texttt{#1}}
}


% Author and Course information
% This Document contains the information about this course.

% Authors of the slides
\author{\href{https://github.com/evilham}{@evilham}\\
\footnotesize{Basierend auf Folien von \href{https://github.com/satkowski}{@satkowski} (November 2016)}}

% Name of the Course
\institute{C\texttt{\#} Kurs}

% Fancy Logo
\titlegraphic{\hfill\includegraphics[height=1.25cm]{../templates/fsr_logo_cropped}}


% Presentation title
\title{Grundlagen von C\texttt{\#} - 2}
\date{\today}


\begin{document}

\maketitle

\begin{frame}{Gliederung}
	\setbeamertemplate{section in toc}[sections numbered]
	\tableofcontents
\end{frame}

% ----------------------- Kontrollstrukturen ------------------------------
\section{Array}
\begin{frame}{Array}
	\begin{itemize}
		\item Sind Aufzählung einer "beliebigen" Anzahl von Variablen gleichen Typs
		\item Können über eine Index aufgerufen werden
		\begin{itemize}
			\item Dieser geht von 0 bis Größe - 1
		\end{itemize}
		\item Können ein- oder mehrdimensional sein
	\end{itemize}
	\textbf{Eindimensionales Array:}\\
	\lstinputlisting{resources/02_grundlagen_2/array.cs}	
\end{frame}

\begin{frame}{Array}
	\textbf{Array von einem Array:}\\
	\lstinputlisting{resources/02_grundlagen_2/array_of_array.cs}
	\textbf{Mehrdimensionales Array:}\\
	\lstinputlisting{resources/02_grundlagen_2/multi_array.cs}	
\end{frame}

% ----------------------- Kontrollstrukturen ------------------------------
\section{Kontrollstrukturen}
\begin{frame}{Kontrollstrukturen}
	\begin{itemize}
		\item Verändern den Ablauf des Programms
		\item Können Verzweigungen oder Schleifen sein
		\item Benötigen bool'sche Ausdrücke als Abfrage
		\begin{itemize}
			\item Kann eine bool-Variable/-Methode sein
			\item oder ein Vergleich
		\end{itemize}
	\end{itemize}
\end{frame}

\subsection{Verzweigungen}
\begin{frame}{Verzweigungen}
	\textbf{if:}\\
	\lstinputlisting{resources/02_grundlagen_2/if.cs}
\end{frame}

\begin{frame}{Verzweigungen}
	\textbf{Tertiäre-Verzweigung:}\\
	\begin{itemize}
		\item Werden genauso genutzt wie eine if-else-Verzweigung
		\item Dabei besteht folgende Syntax:
		\begin{itemize}
			\item \alert{\texttt{<Bedingung> ? <Dann-Pfad> : <Sonst-Pfad>}}
		\end{itemize}
		\item Das \alert{\texttt{?}} steht für \alert{dann}
		\item Das \alert{\texttt{:}} steht für \alert{sonst}
	\end{itemize}		
	\lstinputlisting{resources/02_grundlagen_2/tertiary.cs}
\end{frame}

\begin{frame}{Verzweigungen}
	\textbf{switch case:}\\
	\lstinputlisting{resources/02_grundlagen_2/switch_case.cs}
\end{frame}

\subsection{Schleifen}
\begin{frame}{Schleifen}
	\textbf{while:} \\
	\begin{itemize}
		\item Die Schleife läuft so lange durch deren Körper wie die Bedingung in den Klammern erfüllt ist
	\end{itemize}
	\lstinputlisting{resources/02_grundlagen_2/while.cs}	
\end{frame}

\begin{frame}{Schleifen}
	\textbf{do while:} \\
	\begin{itemize}
		\item Selbe Funktion wie die while-Schleife
		\item Einziger Unterschied ist, dass diese Schleife einmal ihren Körper durchläuft, bevor sie die Bedingung überprüft
	\end{itemize}
	\lstinputlisting{resources/02_grundlagen_2/do_while.cs}	
\end{frame}

\begin{frame}{Schleifen}
	\textbf{for:} \\
	\begin{itemize}
		\item Zählt über eine Zählvariable in gewisser Schrittweite, bis die Bedingung nicht mehr erfüllt ist
		\item Die Schleife hat folgenden Aufbau:
		\begin{itemize}
			\item \alert{\texttt{for( <Zählvariable>; <Bedingung>; <Schrittweite> )}}
		\end{itemize}
	\end{itemize}
	\lstinputlisting{resources/02_grundlagen_2/for.cs}	
\end{frame}

\begin{frame}{Schleifen}
	\textbf{foreach:} \\
	\begin{itemize}
		\item Läuft durch eine Aufzählung an Elementen eines Typen durch 
		\item Weist jeden Durchlauf ein neues Element aus der Aufzählung der Variable zu 
		\item Die Schleife hat folgenden Aufbau:
		\begin{itemize}
			\item \alert{\texttt{foreach( <Datentyp> <Variablenname> in <Aufzählung> )}}
		\end{itemize}
	\end{itemize}
	\lstinputlisting{resources/02_grundlagen_2/foreach.cs}	
\end{frame}

\end{document}
